\part{Classical First-Order Logic}
\label{part:classical-first-order}

Classical first-order logic adds variables ranging over a nonempty domain,
symbols for functions and relations, and quantifiers.  That additional
expressive power makes the syntax of binding and substitution part of the
formal-logical development.  The definitions in this part follow the set-based semantics of
\cite{EbbinghausFlumThomas1994,Hodges1993}; equality is logical identity, not
an unspecified binary relation.

\chapter{First-Order Syntax}
\label{ch:fol-syntax-semantics}

Parallel textbook developments of the material in this part are
\cite{ChiswellHodges2007} and \cite{EbbinghausFlumThomas1994}.

